% AI-written, AI cross-reviewed; no human sign-off; not a claim that RH is proved
% Manuscript v3, heading p. 13; original out/r23_link.tex, SHA-256 742614e691e53727adc35f14c8f1184b1c158a7ad4993d44c3f2d550b014c5bc.
\subsection{From the Selmer sum to the lattice error}\label{r23:link}
Fix one family member $D$, its size datum, and a retained lift $1\le j\le m$.
Write $d_j=\dim Z_{0,j}$ for the reference dimension recorded as
\nolinkurl{dZ0} (\path{DataParams₂.lean:194}), $L=x^\varepsilon$, $b=\exp L$, and $v_\ell=0$ off the primes.
At a prime put $v_\ell=\log\ell\sum_{1\le e<t}\operatorname{shaLen}_{D}(\ell,e)$,
using the rational lengths of this member. Let $S_w$ be the support places
where $\mu(w)\ne\mu_0$, $S$ their residue characteristics, $e_w$ their
ramification indices and $\lambda_w$ the weights \nolinkurl{lwP}.
The geometric error and its non-Selmer allowance are
\[
 E_j=\sum_{w\in S_w}\lambda_w e_w\eta_j(w),\qquad
 A_x=fmbL^2+(\log x)^2N_c.
\]
In these symbols \nolinkurl{LinkN_at} is the implication
$\text{$(v_\ell)_\ell$ summable}\Rightarrow
 E_j\le d_j(CA_x+\sum_\ell v_\ell)$; the family version supplies
one $C$ and one threshold and assumes finiteness of the actual lengths at the
member before quantifying over its lifts.
See \nolinkurl{vvTerm}, \nolinkurl{etaSum}, \nolinkurl{LinkN_at},
\nolinkurl{LinkN} in \path{GlobalTheoremInputs.lean:25--47}.

\begin{formalstatement}[Geometric inputs for the prime sum, inside \decid{D427}]
At a parameter choice $P$ and a scale satisfying
\nolinkurl{Large}$(x,\varepsilon,c)$, assume \nolinkurl{admitted}$(L_{\rm data},P,c)$.
Under finiteness of the lengths this supplies a table with the local
conditions \nolinkurl{AdmT}, including $c>0$ and its polynomial parameter
bounds. Put $H=|\mathrm{Cl}(E)|$, $M=\texttt{restRowsFiniteGroupCard}$,
and let $\lambda_\ell^{G}$ be the table's allowance \nolinkurl{lamG}.
For $\ell\in S$, the contribution over $\ell$ is bounded by $d_j\log\ell$ times
\[
 \nu_\ell+2/\ell^2+\lambda_\ell^{G}+\mathrm{VVr}(\ell)
       +\operatorname{ord}_\ell H+\operatorname{ord}_\ell M,
 \qquad \log\ell\,\mathrm{VVr}(\ell)=v_\ell .
\]
The local inequalities used here are \nolinkurl{etaAt_le_g} on the designated
set and \nolinkurl{etaAt_le} with \nolinkurl{loc_le_lamB} on its complement
(\nolinkurl{TwistedPlaces.etaSum_le}, \path{MemPrimesOverOfMem.lean:558--632}).
For $X=fmbL^2$ and $U=x^{c+2}\exp(x^{\varepsilon/4})$, the following
are the table estimates under these same size hypotheses. In the displayed
order the sums of the five nonnegative envelopes for
$\log\ell\,\lambda_\ell^G$ are bounded by
\[
 2T,\qquad 8cX,\qquad 4cX,\qquad 1000(c+1)^3X,\qquad 1164000T,
 \qquad T=(c+1)^2U^{18}.
\]
The envelopes cover the designated primes, the host's small primes or $3$,
$q,d$, ordinary primes $\ell\ne2,3$ below or at the cutoff, and exceptional
primes above it. The remaining large ordinary primes contribute zero.
Also $6+\nu_{\rm budget}+\log H+\log M\le20T$ and $4\cdot10^6T\le b$.
See \nolinkurl{sum_g}, \nolinkurl{sum_23}, \nolinkurl{sum_qd},
\nolinkurl{sum_small}, \nolinkurl{sum_exc}, \nolinkurl{glob_le}, \nolinkurl{big_le_b},
\nolinkurl{point_bound} in \path{ExportBoundOfAdmissible.lean:18--19},
\path{ExportBoundOfAdmissible.lean:501--502,611--613,651--652},
\path{ExportBoundOfAdmissible.lean:733--735,798--799,825--826,893--996}.
\end{formalstatement}

\begin{proposition}[The local estimate and the family error]
The preceding geometric inputs give the pointwise linking estimate with
$C=2000(|c|+1)^3$ under finiteness and summability. On an eventual
family with $\varepsilon>0$, \nolinkurl{HypV}, \nolinkurl{LinkN} with constant $C$,
and \nolinkurl{SmallLevelRankEq}, every retained lift satisfies, on a common tail,
\[
 E_j\le D_0 (|C|+1)x^\varepsilon(B_0+N_c).
\]
\end{proposition}
\begin{proof}
Sort the filtered places by residue characteristic, counting each place once.
Sum the local bound. Since $H,M$ are positive integers, each truncated
prime factorization contributes at most $\log H$, respectively $\log M$.
The elementary sums in the source give
$\sum_{\ell\in S}\log\ell(\nu_\ell+2/\ell^2)\le6+\nu_{\rm budget}$.
Nonnegativity and summability give
$\sum_{\ell\in S}\log\ell\,\mathrm{VVr}(\ell)\le\sum_\ell v_\ell$.
For the remaining terms, if the prime meets multiple tests, the first
applicable branch of the table is bounded by the sum of their nonnegative
envelopes. Thus the total geometric part inside the factor $d_j$ is at most
\[
 (20+2+1164000)T + (12c+1000(c+1)^3)X\le2000(c+1)^3X.
\]
Indeed the first term is at most $b\le X$ since $f,m,L\ge1$, and
$1+12c\le1000(c+1)^3$. As $N_c\ge0$, $X\le A_x$.
Multiplication by the nonnegative dimension and addition of the Selmer part
give the claimed estimate. The same $c$ all along an admissible family
gives the same $C$ after taking the maximum of the admissibility, largeness
and parameter thresholds. This is the argument of
\nolinkurl{TwistedPlaces.etaSum_le} (\path{MemPrimesOverOfMem.lean:606--647}),
\nolinkurl{export_bound_of_admissible} and \nolinkurl{stmtU_holds}
(\path{ExportBoundOfAdmissible.lean:1023--1053,1141--1160}).

For the family inequality, apply the universal field bound \nolinkurl{VVle}
to the member's own closure and local data. Finiteness allows one to take
the natural value of each extended length, and the record equality
\nolinkurl{sha_eq} identifies it with the rational summand above. Each
$v_\ell$ is nonnegative. Filtering the nonprimes from an arbitrary finite
partial sum bounds it by $\sqrt b$; the supremum of these sums therefore
exists and is at most $\sqrt b$. This proves simultaneously the finiteness,
summability and total bound needed by the family link
(\nolinkurl{UpperTheory.Top.vvRecord}, \path{DivinNOfHypV.lean:17--57}).
Choose the maximum of the three family thresholds, the parameter threshold
plus one, and a threshold beyond which $(\log x)^2\le x^\varepsilon$.
Such a threshold exists since $\log x/x^{\varepsilon/2}\to0$.
Write $y=x^\varepsilon\ge1$. Now $d_j=D_0$, $B_0=fmby$,
$\sqrt b\le b\le B_0$ and $A_x\le y(B_0+N_c)$. Hence
\[
 E_j\le D_0\{ |C|y(B_0+N_c)+B_0\}
       \le D_0(|C|+1)y(B_0+N_c).
\]
See \nolinkurl{log_sq_le_rpow}, \nolinkurl{divinN_at},
\nolinkurl{divinN_of_hypV} in \path{DivinNOfHypV.lean:59--130}.
\end{proof}

\paragraph{Which assumptions pass to the slopes?}
In the application \decid{D427} supplies the family link on the original
rows: construct a good cut for each member above the threshold (use the
whole cut below it), apply the tied-character and two/three-prime bounds
on the cut family, and transfer back. In that transfer the extra term is at
most $d_j\log M$, with $\log M\le C_1x^{C_1}L\le C_1A_x$ eventually
since $x^{C_1}\le b$ and $f,m,L\ge1$. Thus the original constant can
be taken to be the cut constant plus $C_1$. The cut construction and local
compatibility are separate geometric suppliers (see \decid{D421} and \decid{D423})
(\nolinkurl{linkN_of_twGood_galoisPreconditions}, \path{LinkNFam.lean:35--53};
\nolinkurl{linkN_of_gd}, \path{LinkNOfGd.lean:38--93}).
The transfer to divisibility uses the further input \decid{D481}:
on the chosen design, the tuple and discriminant sizes give
\nolinkurl{LinkDIVt}; covered twists (\nolinkurl{LinkEta}) and the bound just
proved give \nolinkurl{LinkDIV}. The remaining closed links, including the
slope and rank data, enter on this same design to produce a tied chain.
These are independent premises of \decid{D488}
(\nolinkurl{LinkDIVt}, \path{GlobalTheoremInputs.lean:77--95};
\nolinkurl{topDesWin}, \path{DesWin.lean:11--31}).

\begin{formalstatement}[The per-unit lattice implication, \nodeid{P021}]
For a finite record with nonempty finite word sets $W(u_0)$ assume
\nolinkurl{HidaSetting}, \nolinkurl{DIV}, \nolinkurl{EComp5}, \nolinkurl{HComp},
\nolinkurl{Count5}, \nolinkurl{Slopes}, \nolinkurl{QuotientSlopeBoundsWith},
\nolinkurl{LossBounds}, \nolinkurl{SlopeIneq}, \nolinkurl{HHBound} and \nolinkurl{FParams}.
Then $\mathrm{ASK}(u_0)\le\mathrm{Main}(u_0)+\mathrm{ErrX5at}(u_0)$.
The clauses and the error polynomial are
\nolinkurl{PointwiseBoundVariant}, \nolinkurl{ErrX5at}
(\path{QuotientSlopeBoundsWith.lean:15--54,72--75};
\path{MeanBoundSecondForm.lean:111--169}).
\end{formalstatement}
\begin{proof}
Fix $u_0$. Write $\langle\cdot\rangle$ for the mean with weights
$d_u/\sum d_u$; these are positive and sum to one. Put $A=\sum_t a_t$,
$s=\langle\mu_N-h\rangle$, $u_F=\max(\mu_N^{\max}-h)-s$,
and $\overline{Lo}=\langle Lo\rangle$. The lower quotient-slope and division
bounds, added, cancel $F_o$ and give at each word
\[
 \mu_N-h\ge \mu_M+G_a+\mu_{CC}-Ee-Lo-S_{\rm Div}
                  +A\log T_u-E_{\rm Div}.
\]
For clarity, the divisibility and counting expressions used below are
\[
\begin{aligned}
 S_{\rm Div}&=m\sum_t\bigl(a_t\log a_t-a_t-a_t(\gamma+\log(\pi/v)-\ell)\bigr),\\
 E_{\rm Div}&=c_{\rm Div}n_TaL_a+c_{\rm Div}m(B_0+N_c)x^{e_{\rm Div}\varepsilon},\\
 \Lambda&=\sum_{j<m}\frac1f\sum_{\rho<f}\log(w_{j\rho}v/\pi)
             +m\,\mathrm{lg}-c_{\rm Count}(L_a+m^2L_a/k),
\end{aligned}
\]
where $\ell$ here is the real logarithmic-derivative entry of the record and
$\mathrm{lg}$ its common logarithmic factor. See \nolinkurl{Data.logOm},
\nolinkurl{Data.Main} (\path{MeanBoundSecondForm.lean:84--109}),
\nolinkurl{Data.SDiv}, \nolinkurl{Data.EDiv} (\path{RankGapProofs/WeightedMeanData.lean:129--134}),
and \nolinkurl{Lam5} (\path{ChainSlopeWindows/LevelOfCongruence.lean:17--19}).
Here $\mu_M\ge-c_Sn_TaL_a(1+m/k)$,
$Ee\le m(A\,\mathrm{lg}+c_E B_0)$, and $\log T_u\ge\Lambda$ by the
lower counting bound, with $A\ge0$. The main-term identity is
\[
 S_{\rm Div}-A\Lambda+mA\,\mathrm{lg}
   =\mathrm{Main}+A c_{\rm Count}(L_a+m^2L_a/k):
\]
expand the sum of $\log(wv/\pi)$ using $w,v>0$; the constant
$mA\log(v/\pi)$ cancels its occurrence in $S_{\rm Div}$ and
$mA\,\mathrm{lg}$ cancels the term from $\Lambda$. Averaging now gives
\[
 \begin{gathered}
 G_a+\mu_{CC}-\mathrm{Main}-K_0-\overline{Lo}\le s,\\
 K_0=A c_{\rm Count}(L_a+m^2L_a/k)+E_{\rm Div}+mc_EB_0+
               c_S n_T aL_a(1+m/k).
 \end{gathered}
\]
Here and below $c_S=c_{\rm Slope}$.
See \nolinkurl{s_lower5}, \nolinkurl{key_identity5}
(\path{ChainSlopeWindows/LevelOfCongruence.lean:17--35,50--66,105--111}).

Compare the upper quotient slope at $u$ to the lower one at $u'$.
The common $G_a+F_o$ cancels. The two $Ee$ errors cost $2mc_EB_0$.
Height comparison costs $c_H(n_TaL_a+mB_0)$ and moving the $A\log T$
terms costs $A c_{\rm Count}\log k$. Write $V_0$ for the sum of these three
costs. The original slope spread costs $c_Sn_TaL_a(2+m/k)$ and the
upper quotient error costs $c_{\rm Sl1}mB_0$. Averaging over $u'$ and
maximizing over $u$ gives $0\le u_F\le K_3+\overline{Lo}$ with
$K_3=c_Sn_TaL_a(2+m/k)+V_0+c_{\rm Sl1}mB_0$.
Nonnegativity follows since $\mu_N\le\mu_N^{\max}$.
By the upper count for $\log T_u$ and $IO\le Lo_P$, \decid{D570} gives
\[
 \overline{Lo}\le M:=c_{\rm Sl1}Lo_P+
 c_{\rm Sl1}n_Tm(a/k+b')c_{\rm Count}mL+c_{\rm Sl1}X''+1 .
\]
This follows directly by substituting the upper count in the formula
for each $Lo(u)$ and averaging; its coefficient is nonnegative.
See \nolinkurl{Vvar_bound5}, \nolinkurl{u_function_le_mean},
\nolinkurl{weighted_mean_lower_bound_le}
(\path{ChainSlopeWindows/LevelOfCongruence.lean:68--103,162--171}).

Set $X=x^{c_{HH}}$. For $z=u_F>0$ and $0<H'=HH\le Xb$, use
$\log t\le t-1$ at $t=(1+H'/z)/(1+X)$:
\[
 z\log(1+H'/z)\le z\log(1+X)+\frac{H'-zX}{1+X}
 \le(z+b)\log(1+X).
\]
The last inequality uses $X/(1+X)\le\log(1+X)$.
At $z=0$ the source uses zero for the left side. The finite parameter bounds
give $0\le\delta\le\bar\delta<1$, $0\le\eta\le\bar\eta$,
$0\le e\le\bar e\le1/4$, so the height hypothesis permits application of
the slope inequality. It yields
\[
 s\le\mu_{CC}+\log|W(u_0)|+
 \frac{4\log(1+X)(K_3+M+b)+(\bar\eta+4\sqrt{\bar e})Xb}{1-\bar\delta}.
\]
Cancel $\mu_{CC}$ against the lower bound for $s$, insert
$\mathrm{ASK}-G_a\le Lo_P$ and $\overline{Lo}\le M$.
The resulting error is $K_0+\log|W(u_0)|+Lo_P+M$ plus the displayed
fraction, exactly \nolinkurl{ErrX5at}
(\nolinkurl{Ga_bound5}, \nolinkurl{each_unit_bound_statement},
\path{ChainSlopeWindows/LevelOfCongruence.lean:113--160,175--197}).
Finally the top and place bijections, with the matching weight and scalar
identities, take $\mathrm{Main}$ to $\sum_{t,j}\Re\mathrm{MAIN}$.
The fifth field of \nolinkurl{IdentU} independently identifies $\mathrm{ASK}$
with the arithmetic sum (\path{DedekindLFunctions/SplitInertia.lean:67--72}). Subtraction gives
$\operatorname{gap}(\operatorname{tr}u_0)\le\mathrm{ErrX5at}(u_0)$;
the two additional size-window bounds and large-$x$ condition in
\nolinkurl{ch_top} convert that error to the bracket of the previous stage.
See \decid{D569}, \nolinkurl{ch_top_main_eq}, \nolinkurl{ch_top}
(\path{LandauTheorems/ChTop.lean:11--55}).
\end{proof}
